%% ============================================================
%% DeIR-Dual: Dual-Query Instruction Rewriting for Training-Free
%% Instruction-Following Retrieval
%%
%% Target: AAAI 2027
%% Format: AAAI two-column LaTeX
%% ============================================================

\documentclass[letterpaper]{article}
\usepackage{aaai26}
\usepackage{times}
\usepackage{helvet}
\usepackage{courier}
\usepackage[T1]{fontenc}
\usepackage{hyperref}
\usepackage{amsmath}
\usepackage{amssymb}
\usepackage{booktabs}
\usepackage{graphicx}
\usepackage{subcaption}
\usepackage{algorithm}
\usepackage{algpseudocode}
\usepackage{multirow}
\usepackage{enumitem}
\usepackage{xcolor}

% ============================================================
% Custom macros
% ============================================================
\newcommand{\method}{DeIR-Dual}
\newcommand{\qbase}{Q_{\text{base}}}
\newcommand{\qpos}{Q_{\text{pos}}}
\newcommand{\qminus}{Q_{\text{minus}}}
\newcommand{\sbase}{S_{\text{base}}}
\newcommand{\spos}{S_{\text{pos}}}
\newcommand{\sneg}{S_{\text{neg}}}
\newcommand{\sfinal}{S_{\text{final}}}
\newcommand{\todo}[1]{\textcolor{red}{\textbf{[TODO] #1}}}

\hypersetup{pdfinfo={TemplateVersion={2026.1}}}

\setcounter{secnumdepth}{2}

% ============================================================
% Author info (uncomment for final submission)
% ============================================================
% \author{
%   Author One,
%   Author Two,
%   Author Three
% }
% \affiliations{
%   Affiliation 1\\
%   Affiliation 2\\
%   Affiliation 3\\
% }

\begin{document}

\title{DeIR-Dual: Dual-Query Instruction Rewriting for Training-Free Instruction-Following Retrieval}
\author{Anonymous Authors}
\date{}
\maketitle

% ============================================================
% ABSTRACT
% ============================================================
\begin{abstract}
Instruction-following retrieval requires systems to adapt their ranking
when query instructions change. Existing query rewriting methods
(e.g., HyDE, Query2Doc, DeepRetrieval) improve raw retrieval quality
through semantic enhancement but lose instruction sensitivity
(p-MRR $\approx$ 0). We propose \method{}, which decomposes an
instruction into a positive enhancement query ($\qpos$) and a
negative exclusion query ($\qminus$), then applies a training-free
dual-track scoring mechanism with a dynamic semantic threshold
$\tau = \cos(\mathbf{h}_{\text{base}}, \mathbf{h}_{\text{neg}}) + \delta$. Our key findings are:
(1) a first-principles \emph{Scale Alignment} derivation converges
to $\alpha = 1.0$ across 30 independent methods, improving p-MRR by
70.9\% over grid search; (2) \method{} outperforms six baselines by
9--44$\times$ in p-MRR while achieving the highest retrieval quality
(target\_avg) on the FollowIR benchmark; (3) cross-modal validation
across 7+ datasets confirms broad applicability; and (4) we identify
and analyze a \emph{semantic entanglement} failure mode with a
conservative prompting remedy.

Code: \url{https://github.com/...}
\end{abstract}

% ============================================================
% 1. INTRODUCTION
% ============================================================
\section{Introduction}
\label{sec:introduction}

Query rewriting has emerged as a powerful paradigm for improving dense
retrieval, with methods such as HyDE \cite{gao2023hyde}, Query2Doc
\cite{wang2023query2doc}, and DeepRetrieval \cite{li2025deepretrieval}
demonstrating substantial gains in retrieval quality by enriching the
query's semantic coverage.  Some methods even generate multiple query
variants --- RAG-Fusion \cite{raudaschl2024ragfusion} produces diverse
rephrasings and fuses their retrieval results.  However, all existing
rewriting methods share a critical limitation: their rewritten queries
are \textbf{semantically unidirectional} --- they all point toward
``what to retrieve'' (semantic enrichment) but none creates an
independent channel for ``what to exclude.''  Even RAG-Fusion's
multiple variants are all positive rephrasings of the same intent; none
explicitly captures the negative exclusion signal.  This makes current
rewriting fundamentally \textbf{insensitive to instruction changes}:
when the exclusion constraint in an instruction is modified (e.g.,
``exclude repair-only articles'' $\to$ ``exclude funding-related
articles''), the rewritten queries barely change their ranking, yielding
near-zero instruction sensitivity (p-MRR $\approx$ 0)
\cite{followir2024}.  In short, current rewriting can express
\emph{where to look} but cannot independently express \emph{where not
to look}.

Why do existing rewriting methods fail at instruction discrimination?
The root cause is that real-world retrieval instructions contain
\textbf{three distinct types of signals} --- background context,
positive refinement, and negative exclusion \cite{followir2024} ---
and all current rewriting methods, whether single-query or multi-query,
address only the first two.  The negative exclusion signal --- the
directive about what documents are \emph{not} relevant --- has no
independent channel in any existing rewriting framework.  Collapsing
it together with positive signals causes inevitable interference: the
exclusion signal is diluted by the dominant positive and background
signals, and there is no mechanism to selectively suppress documents
that match the exclusion constraint.  This tension between
\textbf{semantic coverage} (retrieving more relevant content) and
\textbf{instruction discrimination} (adapting the ranking to
instruction changes) lies at the heart of the instruction-following
retrieval problem, and current rewriting methods address only the
former.

Prior work has attempted to address instruction discrimination from
other angles, but each falls short.  Supervised instruction tuning ---
Promptriever \cite{weller2024promptriever} and Dual-View Training
\cite{dualview2026} --- improves instruction following through
fine-tuning on annotated data, but the annotation cost is high and
generalization to unseen exclusion types remains fragile.  Negation-aware
methods such as DEO \cite{deo2026} decompose queries into
positive and negative components at the representation layer, but rely
on linear vector-space operations ($Q_{\text{base}} - Q_{\text{neg}}$)
that distort the positive intent when moving away from the negation
direction.  Critically, all three paradigms --- semantic rewriting,
supervised tuning, and representation-layer decomposition --- share a
common assumption: that instruction signals must be resolved at the
\textbf{representation layer}, by producing a single (or linearly
combined) query embedding.

We challenge this assumption from within the rewriting paradigm itself.
The key insight is that the failure of current rewriting is not due to
poor rewrite quality, but due to \textbf{structural incompleteness}:
no amount of positive-only rewriting, whether single or multiple
variants, can create an independent exclusion channel.  Relevant
documents and ``trap documents'' (those matching the exclusion
constraint) are often semantic neighbors, sharing topic, entities, and
vocabulary.  In such entangled regions, no collection of positive
rephrasings can cleanly separate them.  The solution is not to rewrite
\emph{more}, but to rewrite \emph{structurally}: decompose the
instruction into independent channels --- including an explicit
exclusion channel --- and resolve their signals at the \textbf{decision
layer}, where non-linear logic can selectively suppress trap documents
without distorting the embedding geometry.

We propose \method{} (\textbf{De}composition-based
\textbf{I}nstruction-\textbf{R}ewriting with \textbf{Dual}-track
scoring), a training-free framework that operates entirely at the
scoring layer.  \method{} first prompts an LLM to decompose an
instruction into a positive query ($\qpos$) and a negative query
($\qminus$), which are encoded alongside the base query ($\qbase$)
through a frozen dense retriever.  These three channels feed into a
dual-track scoring mechanism: a \textbf{penalty track} with a dynamic
semantic threshold $\tau = \cos(\mathbf{h}_{\text{base}},
\mathbf{h}_{\text{neg}})$ that anchors the exclusion boundary to the
query-instruction affinity, and a \textbf{reward track} gated by a
conditional Safety Gate that severs the enhancement signal when a
document triggers the exclusion constraint.  This non-linear coupling
--- which we show implements a logical-NOT in embedding space --- is
what linear vector methods cannot replicate, and it comes at zero
training cost: \method{} requires no fine-tuning, no embedding
modification, and no annotated instruction pairs.

Our experiments reveal four findings.  First, scoring parameters need
not be empirically tuned: a \textbf{Scale Alignment Principle} derived
from the intrinsic statistics of the embedding space yields
$\alpha = 1.0$ across 30 independent derivation strategies, improving
p-MRR by 70.9\% over grid search while eliminating test-set leakage.
Second, on the FollowIR benchmark \cite{followir2024}, \method{}
outperforms six baselines by 9--44$\times$ in p-MRR while achieving
the highest retrieval quality (target\_avg) --- despite never being
trained on an instruction-annotated pair.  Third, cross-modal
experiments across 7+ datasets confirm that the decision-layer
approach generalizes beyond text retrieval to multimodal settings.
Fourth, we identify \textbf{semantic entanglement} --- where the
exclusion vector overlaps substantially with the base intent --- as a
fundamental failure mode, and characterize when and why the dual-track
mechanism breaks down.

In summary, our contributions are: \textbf{(i)} \method{}, a
training-free decision-layer framework that shifts
instruction-following from a training problem to a scoring problem,
achieving the highest p-MRR and retrieval quality on FollowIR without
model fine-tuning; \textbf{(ii)} the Scale Alignment Principle, a
principled calibration framework that derives scoring parameters from
the retriever's own score distributions rather than from test-set
grid search; and \textbf{(iii)} an empirical characterization of
semantic entanglement, the key failure mode limiting training-free
instruction-following retrieval.

% ============================================================
% 2. RELATED WORK
% ============================================================
\section{Related Work}
\label{sec:related-work}

\subsection{Query Rewriting and Semantic Enhancement}

Query rewriting aims to bridge the vocabulary gap between user queries
and document collections.  Early approaches include HyDE
\cite{gao2023hyde}, which generates hypothetical documents,
Query2Doc \cite{wang2023query2doc}, which concatenates queries with
pseudo-documents, and RAG-Fusion \cite{raudaschl2024ragfusion}, which
fuses multiple LLM-generated variants.  RAG-QR
\cite{ma2023ragqr} learns a rewriter via reinforcement learning, while
DeepRetrieval \cite{li2025deepretrieval} fine-tunes a language model
with RL to optimize retrieval recall.  Recent works such as
FollowTable \cite{followtable2026} extend instruction-following
retrieval to structured data domains.  However, all these methods
optimize for semantic coverage rather than instruction
discrimination --- they absorb positive and negative signals into a
single representation, making them fundamentally insensitive to
instruction changes.  This is confirmed empirically: semantic
enhancement methods achieve near-zero instruction sensitivity
(p-MRR $\approx 0$) despite strong baseline retrieval quality.

\subsection{Instruction-Following Retrieval}

The success of instruction-following language models such as
InstructGPT \cite{instructgpt2022}, FLAN \cite{flan2022}, and T0
\cite{t02022} has prompted the NLP community to explore instruction
tuning across diverse domains \cite{llama2023,groeneveld2024}.  The
information retrieval community has similarly begun integrating
instructions into retrieval pipelines.  INSTRUCTOR
\cite{instructor2022} and TART \cite{tart2022} were among the earliest
to train retrievers with task-specific instructions, though these
instructions are typically short task prefixes rather than query-level
constraints.  Recent efforts scale up model size and instruction
coverage: E5-Mistral \cite{e5mistral2023}, GritLM \cite{gritlm2024},
and BGE \cite{bge2023} use LLaMA or Mistral backbones to support
richer instruction sets, while C-Pack \cite{bge2023} provides
packaged resources for multilingual embedding training.  Despite this
activity, existing instruction-aware retrievers evaluate on standard
benchmark suites such as MTEB \cite{mteb2022} and BEIR
\cite{beir2021} that contain no per-query instructions
\cite{followir2024}.  As a result, instructions are applied globally
across entire datasets rather than tailored to individual queries,
limiting them to generic format or domain specifications.  The
FollowIR benchmark \cite{followir2024} addresses this gap by
introducing per-query instructions with the original/changed query
pair protocol and the p-MRR metric.  Promptriever
\cite{weller2024promptriever} trains retrievers to follow natural
language instructions, InF-IR \cite{zhuang2025infir} releases a
large-scale instruction-following training corpus and model, while
Dual-View Training
\cite{dualview2026} synthesizes training data via polarity reversal.
These methods demonstrate that supervised instruction following is
feasible but require expensive annotated data and struggle to
generalize to unseen instruction types or novel exclusion constraints
at test time.

\subsection{Negation-Aware Retrieval}

Handling negation in retrieval has received growing attention across
both text and multimodal domains.  The DEO framework
\cite{deo2024,deo2026} learns separate representations for positive
and negative instruction components in text retrieval, with the latest
work \cite{deo2026} proposing a training-free approach that operates at
the representation layer by optimizing query embeddings directly.
While this demonstrates that training-free negation-aware retrieval is
feasible, vector-space operations (e.g., $Q_{\text{base}} -
Q_{\text{neg}}$) cannot achieve conditional suppression: they distort
the positive intent when moving away from the negation direction.  In
the multimodal domain, NegCLIP \cite{negclip2022} fine-tunes CLIP on
negation-augmented data to distinguish positive from negative
descriptions, while COCO-Neg \cite{coconeg2024} reveals that standard
models catastrophically fail under caption polarity inversion.

Our approach differs fundamentally from all prior work by operating at
the \textbf{decision layer} rather than the representation layer.
Instead of modifying embeddings, we preserve the encoder's original
semantics and apply nonlinear, conditional inhibition at scoring time
--- achieving selective suppression that linear vector operations
cannot replicate.

% ============================================================
% 3. METHOD
% ============================================================
\section{Method: DeIR-Dual}
\label{sec:method}

% Section narrative: We first formalize the structural blind spot
% in single-vector retrieval, then introduce the dual-query
% decomposition as the key insight, present the scoring mechanism
% with its nonlinear Safety Gate, and finally describe the
% Scale Alignment principle for parameter derivation.

\subsection{Problem Formulation}
\label{sec:problem-formulation}

Dense retrievers encode a query into a single dense vector
$\mathbf{v}_{q,I} \in \mathbb{R}^d$ and retrieve documents by descending
order of $\cos(\mathbf{v}_{q,I}, \mathbf{v}_d)$.  As shown in FollowIR
\cite{followir2024}, real-world retrieval instructions contain \textbf{three
distinct types of signals} that are structurally incompatible with
single-vector encoding:
%
\begin{enumerate}
  \item \textbf{Background context} (\emph{less directly-relevant background
    info}): peripheral information that helps the system understand the
    user's information need but does not directly determine relevance.
  \item \textbf{Positive refinement} (\emph{more specific details about what
    is relevant}): fine-grained constraints that narrow the relevance
    definition (e.g., ``must contain information about formation'').
  \item \textbf{Negative exclusion} (\emph{directives about what documents
    are not relevant, using negation}): explicit constraints that mark
    certain documents as irrelevant (e.g., ``documents limited to the
    problems of the telescope are irrelevant'').
\end{enumerate}
%
When all three signals are collapsed into a single vector, they
inevitably interfere: the background context dilutes the positive
refinement, and the negative exclusion has no independent channel to
express ``where not to look.''  This is the \textbf{structural blind spot}
of single-vector retrieval --- it can express semantic direction but not
semantic boundaries.

Instruction sensitivity is measured by p-MRR
\cite{followir2024}: how much a system's ranking changes when the
instruction changes, with zero or negative values indicating failure
to adapt.

The core challenge is therefore: \textbf{how can we create independent
channels for all three instruction signals --- background context,
positive refinement, and negative exclusion --- without modifying the
encoder's embedding space?}

\subsection{Triple-Query Decomposition}
\label{sec:dual-query-generation}

The key insight of \method{} is to create \textbf{three independent
channels} in the scoring layer, each handling one class of instruction
signal from FollowIR \cite{followir2024}:

\begin{enumerate}
  \item \textbf{$\qbase$ (Base Query)}: the original query text
    concatenated with the full instruction.  This channel preserves
    the \textbf{background context} --- all contextual information that
    helps frame the information need, including the user's intent,
    domain, and peripheral constraints.
  \item \textbf{$\qpos$ (Positive Refinement)}: an LLM-generated query
    that extracts and strengthens the \textbf{positive refinement signals}
    from the instruction --- the ``more specific details about what is
    relevant'' (e.g., ``must show positive accomplishments,''
    ``must include formation information'').  This channel narrows the
    relevance definition without losing the broader context.
  \item \textbf{$\qminus$ (Negative Exclusion)}: an LLM-generated query
    that explicitly captures the \textbf{negative exclusion signals} ---
    the ``directives about what documents are not relevant, using
    negation'' (e.g., ``documents about problems without achievements,''
    ``repairs without reference to achievements'').  This channel
    creates an independent ``do not retrieve'' signal.
\end{enumerate}

The three queries are independently encoded by the same dense retriever
$f(\cdot)$, producing three query embeddings
$\mathbf{h}_{\text{base}} = f(\qbase)$,
$\mathbf{h}_{\text{pos}} = f(\qpos)$,
$\mathbf{h}_{\text{neg}} = f(\qminus)$.
Document embeddings $\mathbf{v}_d = f(d)$ are pre-computed once and reused.

\textbf{Why three channels, not two?}  Prior work (e.g., DEO
\cite{deo2024}) decomposes instructions into positive and negative
components only.  However, this overlooks the \textbf{background context}
signal --- the peripheral information that is not directly relevant but
frames the user's intent.  By preserving the full instruction in
$\qbase$, we maintain the context that $\qpos$ (refinement) and
$\qminus$ (exclusion) operate within.  Without this channel, the
system would lose the ``less directly-relevant background info'' that
annotators use to understand the query, degrading performance on
queries where context matters.

\subsection{Dual-Track Scoring and Non-Linear Coupling}
\label{sec:scoring-formula}

Given the three query encodings, we compute cosine similarities:
$\sbase = \cos(\mathbf{h}_{\text{base}}, \mathbf{v}_d)$,
$\spos = \cos(\mathbf{h}_{\text{pos}}, \mathbf{v}_d)$,
$\sneg = \cos(\mathbf{h}_{\text{neg}}, \mathbf{v}_d)$.
Unlike simple linear query expansion, our framework explicitly models the
interaction between intent enhancement and exclusion through a
\textbf{non-linear coupling logic}.  We organize the mechanism into
two interacting tracks, whose formulas are introduced in context
below and summarized in Eqs.~\ref{eq:tau}--\ref{eq:final}.

Importantly, this dual-track scoring requires \textbf{no parameter
updates} to the base retriever $f(\cdot)$: all operations occur at
the scoring layer, making \method{} a 100\% training-free inference
engine.

\subsubsection{The Penalty Track: Adaptive Exclusion Boundary}

The primary challenge in exclusion-aware retrieval is distinguishing
between genuine constraint violation and incidental topical overlap.
The penalty track reshapes the decision surface around the exclusion
signal through three coordinated components.

\textbf{Semantic Noise Floor ($\tau$).}
Recall that $\mathbf{h}_{\text{base}}$ encodes the full background
context while $\mathbf{h}_{\text{neg}}$ encodes the pure exclusion
signal.  A fixed penalty threshold fails because the
\textbf{topological overlap} between these two vectors varies
dramatically across queries.  We define a \textbf{dynamic semantic
threshold} that measures precisely this overlap:
%
\begin{equation}
  \tau = \cos(\mathbf{h}_{\text{base}}, \mathbf{h}_{\text{neg}}) + \delta.
  \label{eq:tau}
\end{equation}
%
When the two directions are nearly orthogonal (e.g., ``ISO
standards,'' ``exclude IEEE''), $\tau \ll 0$ and the penalty deploys
aggressively, filtering trap documents with confidence.  When they are
semantically entangled (e.g., ``Hubble achievements,'' ``exclude
repair-only articles''), $\tau$ rises toward $1.0$, automatically
raising the penalty trigger bar and protecting documents that merely
share vocabulary with the exclusion.  This dynamic anchoring is
precisely why preserving $\mathbf{h}_{\text{base}}$ as an independent
channel is essential: without it, we cannot measure the topological
overlap that determines whether an exclusion should trigger aggressive
or conservative filtering.  The offset $\delta$ is set to $0.0$ under
the Zero-Bias Criterion (\S\ref{sec:delta-derivation}).

\textbf{Relative Saliency Gating ($\text{gap}_w$).}
A document should only be penalized when its exclusion signal
\emph{dominates} its base relevance --- i.e., the document is
genuinely more about ``what to exclude'' than ``what to retrieve.''
We formalize this through a saliency gate:
%
\begin{equation}
  \text{gap}_w = \sigma\big((\sneg - \sbase) \cdot T_{\text{gap}}\big).
  \label{eq:gap}
\end{equation}
%
When $\sneg > \sbase$, the gate opens ($\text{gap}_w \to 1$) and the
penalty activates.  When $\sbase$ dominates ($\sbase \gg \sneg$), the
gate closes ($\text{gap}_w \to 0$), protecting topically relevant
documents from collateral damage.  This \textbf{context-aware} design
preserves topical integrity: a document that is strongly relevant on
the base channel is shielded from the exclusion mechanism, even if it
happens to have a modest exclusion score.

\textbf{Smooth Barrier (Softplus with Saturation).}
We apply the penalty through a Softplus-transformed barrier, capped
by a saturation factor $\rho$:
%
\begin{equation}
  \text{penalty} = \min\!\big(\alpha \cdot \text{Softplus}(\sneg - \tau)
    \cdot \text{gap}_w,\; \sbase \cdot \rho\big).
  \label{eq:penalty}
\end{equation}
%
The Softplus transformation (rather than a hard ReLU cutoff) creates a
\textbf{continuous penalty landscape} that prevents abrupt
rank-swapping anomalies in the retrieval manifold.  The saturation cap
$\rho = 0.5$ prevents the \textbf{over-killing effect}: a single
exclusion term should not be able to zero out a document's total
relevance.  The penalty scale $\alpha$ is derived from first
principles in \S\ref{sec:parameter-derivation}.

\subsubsection{The Reward Track: Conditional Enhancement}

A naive addition of an enhancement score $\spos$ leads to
\textbf{reward leakage}: trap documents that match the topic but
violate the instruction still receive a boost, defeating the purpose
of exclusion.  We resolve this through conditional gating.

\textbf{Safety Gate.}
We implement a non-linear gate that monitors the penalty track and
suppresses the reward channel when a document enters the exclusion
zone:
%
\begin{equation}
  \text{safety} = 1 - \sigma\big((\sneg - \tau) \cdot T_{\text{safety}}\big).
  \label{eq:safety}
\end{equation}
%
The gate is near-binary ($T_{\text{safety}} = 20$): for documents
safely below threshold ($\sneg \ll \tau$), $\text{safety} \approx 1$
and the full reward applies; for documents that trigger the exclusion
zone ($\sneg \gg \tau$), $\text{safety} \approx 0$ and the reward is
\textbf{completely cut off}.

\textbf{Coupled Synthesis.}
The final score is produced by coupling the two tracks:
%
\begin{equation}
  \sfinal = \sbase + \beta \cdot \spos \cdot \text{safety} - \text{penalty}.
  \label{eq:final}
\end{equation}
%
This configuration implements a \textbf{Logical-NOT} operation within
the embedding space: reward is conditionally granted only if the
document resides in the ``safe'' zone.  The non-linear synergy between
Eqs.~\ref{eq:safety} and \ref{eq:final} cannot be replicated by linear
vector-space operations ($Q_{\text{base}} - Q_{\text{neg}}$), which
inevitably distort the positive intent when moving away from the
negation direction.  The enhancement scale $\beta$ is derived in
\S\ref{sec:parameter-derivation}.

\textbf{Complete Formula.}
Substituting Eqs.~\ref{eq:tau}--\ref{eq:safety} into
Eq.~\ref{eq:final}, the full scoring function of \method{} is:
%
\begin{equation}
  \boxed{
  \sfinal = \sbase + \beta\,\spos\cdot\text{safety} - \min\!\big(\alpha\,\text{Softplus}(\sneg-\tau)\cdot\text{gap}_w,\; \sbase\cdot\rho\big)
  }
  \label{eq:complete}
\end{equation}
%
where $\tau = \cos(\mathbf{h}_{\text{base}}, \mathbf{h}_{\text{neg}}) + \delta$,
$T_{\text{safety}} = 20$, $T_{\text{gap}} = 10$, $\rho = 0.5$,
and the three learnable-free parameters $\alpha$, $\beta$, $\delta$
are derived from first principles in
\S\ref{sec:parameter-derivation}.

\begin{figure}[t]
  \centering
  % \includegraphics[width=\linewidth]{figures/method_pipeline.pdf}
  \fbox{\parbox{\linewidth}{\centering\vspace{3cm}%
    \textbf{Figure 1: \method{} Pipeline}\\
    \small (placeholder: to be replaced with pipeline diagram)\vspace{1cm}}}
  \caption{Overview of \method{}.  A query and its instruction are
    decomposed by an LLM into three queries ($\qbase$, $\qpos$,
    $\qminus$), independently encoded, and scored via the nonlinear
    decision mechanism in Eq. \ref{eq:final}.  The Safety Gate
    (Eq. \ref{eq:safety}) conditionally suppresses the enhancement
    signal when a document triggers the exclusion threshold $\tau$,
    enabling what linear vector operations cannot: selective,
    conditional suppression that penalizes trap documents while
    preserving relevant document rankings.}
  \label{fig:pipeline}
\end{figure}

% ============================================================
% 4. FIRST-PRINCIPLES PARAMETER DERIVATION
% ============================================================
\section{Scale Alignment: From Empirical Search to Statistical Calibration}
\label{sec:parameter-derivation}

The scoring formula in Eq.~\ref{eq:final} contains three parameters:
$\alpha$ (penalty magnitude), $\beta$ (enhancement magnitude), and
$\delta$ (threshold offset).  Rather than grid-searching these on the
test set --- which risks overfitting and provides no insight into why a
particular value is optimal --- we derive them from the intrinsic
statistical properties of the score distributions.  The core
contribution of this section is a principled calibration framework that
transforms parameter selection from an empirical search into a
statistical derivation.

\subsection{The Symmetry of Reward and Penalty ($\alpha$ and $\beta$)}
\label{sec:scale-alignment}

To ensure the instruction-following mechanism is both effective and
robust across different encoders, we move beyond empirical tuning and
propose the \textbf{Scale Alignment Principle (SAP)}: the
``instructional energy'' (either positive or negative) should be
commensurate with the ``topic energy'' (the base retrieval score)
within the embedding space.  We define
$E_{\text{base}}(d) \equiv S_{\text{base}}(d)$ as the document's
\textbf{signal intensity} on the base channel, and analogously
$E_{\text{req}}(d)$ and $E_{\text{neg}}(d)$ for the refinement and
exclusion channels.  SAP treats $\alpha$ and $\beta$ as scaling
factors that translate heterogeneous instructional signals into a
common probabilistic scale, ensuring \textbf{magnitude symmetry}
between the base retriever and our dual-track modifications.

This calibration philosophy has strong precedents.  In classification,
temperature scaling \cite{guo2017temperature} and Platt scaling
\cite{platt1999probabilistic} derive a single post-hoc parameter from
the statistics of logit distributions, without retraining the model.
In information retrieval, the Divergence from Randomness framework
\cite{amati2002dfr} derives scoring functions from first principles of
probability theory rather than empirical tuning.  SAP extends this
tradition to instruction-following retrieval: instead of calibrating
confidence estimates, we calibrate the relative magnitude of
instructional correction signals against the base retriever's scores,
transforming three heterogeneous similarity scores into a
well-calibrated ranking signal.

\subsubsection{Penalty Calibration ($\alpha$): Neutralizing Topical Bias}

We define ``at-risk'' documents as those where
$S_{\text{neg}} > \tau$ (i.e., documents that match the exclusion
constraint).  For these documents, the penalty must be sufficient to
\textbf{neutralize the existing topical momentum} --- the force to
exclude should equal the force that originally retrieved the document.
Applying SAP to the penalty channel yields:
%
\begin{equation}
  \alpha_{\text{SAP}} =
    \frac{\mathbb{E}[S_{\text{base}} \mid \text{at-risk}]}
         {\mathbb{E}[\text{Softplus}(S_{\text{neg}} - \tau) \mid
         \text{at-risk}]}.
  \label{eq:alpha}
\end{equation}
%
Eq.~\ref{eq:alpha} requires no hyper-parameter search, no reward
function, and no heuristic assumptions --- it simply calibrates the
penalty to match the base signal.

To verify that $\alpha_{\text{SAP}}$ is not an artifact of a single
derivation strategy, we evaluated \textbf{30 independent derivation
methods} spanning six theoretical families
(Table~\ref{tab:alpha-derivation}; full list in Appendix
\ref{sec:alpha-appendix}):
%
\begin{itemize}
  \item \textbf{Group A (Primary: Scale Alignment)}: expectation ratio and
    percentile-matching methods.  $\alpha = 1.000$--$1.030$.
  \item \textbf{Groups B--E (Robustness Checks)}: score resolution
    statistics, distribution separation (KL divergence), ranking-margin
    analysis, and physics-informed half-life models.  These produce lower
    $\alpha$ values ($0.05$--$0.50$) that are overly conservative.
  \item \textbf{Group F (Boundary Exploration)}: entropy-based resolution,
    kurtosis scaling, and Bayesian posterior estimates.  These yield
    $\alpha$ values up to $6.15$ --- mathematically valid but practically
    too aggressive, violating a \emph{Minimal Intervention} principle.
\end{itemize}
%
Across all 30 strategies, the median $\alpha = 0.92$ with $\sigma =
0.15$, and Group~A methods exhibit a \textbf{surprising consensus}
around $\alpha \approx 1.0$ under the standard setting ($\delta =
0.0$).  This cross-method agreement suggests that $\alpha = 1.0$ is
not merely an empirical hyper-parameter, but a \textbf{fundamental
scaling constant} that maintains magnitude symmetry between semantic
relevance and instructional exclusion in dense embedding spaces.

\subsubsection{Enhancement Calibration ($\beta$): The Information Premium}

For ``safe'' documents ($S_{\text{neg}} \leq \tau$), the enhancement
magnitude should match the base score magnitude.  Applying SAP to the
reward channel yields:
%
\begin{equation}
  \beta = \frac{\mathbb{E}[S_{\text{base}} \mid \text{safe}]}
    {\mathbb{E}[S_{\text{req}} \cdot \text{safety} \mid \text{safe}]}.
  \label{eq:beta}
\end{equation}
%
From training-set statistics (855 queries, 878 positives, 12,825
negatives), this yields $\beta = 1.926$.  The fact that $\beta > 1$
reveals a fundamental \textbf{Information Asymmetry}: while $\alpha$
merely \emph{neutralizes} existing topical bias (working against the
base score), $\beta$ must \emph{inject} new semantic information from
$\qpos$ that is not present in $\qbase$.  This \textbf{Information
Premium} --- the extra multiplier beyond $1.0$ --- quantifies the
cost of ``promoting'' documents that satisfy specific instructional
constraints above the generic topical background.  In other words,
$\alpha$ is a damping factor while $\beta$ is an amplification factor,
and their asymmetry ($\beta > 1$ vs.\ $\alpha \approx 1$) reflects
the different nature of the two instructional channels.

\begin{table}[t]
\centering
\scriptsize
\begin{tabular}{lcccc}
\toprule
\textbf{Group} & \textbf{Method} & $\alpha$ & \textbf{Intuition} & \textbf{Status} \\
\midrule
A & Scale Alignment & 1.000 & Expectation ratio & \textbf{Optimal} \\
A & Percentile-50 & 1.000 & Median alignment & Consistent \\
A & Percentile-75 & 1.030 & 75th-percentile alignment & Near-optimal \\
B & Score Variance & 0.052 & Encoder resolution & Too low \\
C & KL Minimization & 0.120 & Distribution match & Too low \\
D & NDCG Margin & 0.310 & Ranking margin & Suboptimal \\
E & Half-Life & 0.500 & 50\% decay & Conservative \\
F & Score Entropy & 6.150 & Shannon resolution & Too high \\
\bottomrule
\end{tabular}
	\caption{$\alpha$ derivation across 30 candidate strategies
	(representative subset shown; full list in Appendix
	\ref{sec:alpha-appendix}).  Group~A methods exhibit a striking
	consensus around $\alpha \approx 1.0$.  Groups B--E are
	conservative (under-penalize), while Group~F explores boundary
	cases that violate the Minimal Intervention principle.}
\label{tab:alpha-derivation}
\end{table}

\subsection{The Zero-Bias Threshold ($\delta$)}
\label{sec:delta-derivation}

Rather than treating $\delta$ as a free parameter, we derive it from
the intrinsic noise characteristic of the exclusion channel.  We
define:
%
\begin{equation}
  \delta = k \cdot \sigma_{\text{neg}},
  \label{eq:delta}
\end{equation}
%
where $\sigma_{\text{neg}} = \text{Std}(S_{\text{neg}})$ is the
standard deviation of the exclusion scores across the training corpus
--- measuring the channel's random fluctuation amplitude --- and $k$
is a confidence multiplier.

This formulation admits two principled choices for $k$, grounded in
distinct statistical criteria:
%
\begin{itemize}
  \item \textbf{$k = 0.0$ (Neyman--Pearson threshold)}: the dynamic
    anchor $\tau = \cos(\mathbf{h}_{\text{base}}, \mathbf{h}_{\text{neg}})$
    already captures the per-query semantic noise floor.  Adding zero
    noise margin maximizes the at-risk population and hence instruction
    sensitivity.  This is our primary setting.
  \item \textbf{$k = 2.0$ (95\% confidence margin)}: adds two standard
    deviations of the exclusion score as a noise buffer.  While
    statistically conservative, this raises the penalty trigger bar
    and shrinks the at-risk population, reducing p-MRR.
\end{itemize}
%
From training-set statistics, $\sigma_{\text{neg}} = 0.087$, yielding
$\delta = 0.0$ under the Neyman--Pearson criterion ($k = 0.0$).  The
fact that $\delta$ is \emph{computed} from Eq.~\ref{eq:delta} rather
than asserted distinguishes our approach from heuristic threshold
tuning.

More fundamentally, \textbf{any $\delta \neq 0$ would break the SAP
equilibrium}.  Since $\alpha$ and $\beta$ are calibrated under the
$\delta = 0.0$ baseline (Eq.~\ref{eq:delta} with $k = 0.0$),
introducing a non-zero offset shifts the at-risk population and
unbalances the symmetry between penalty magnitude and base score
magnitude established by Eq.~\ref{eq:alpha}.  Empirically, $k > 0$
restricts the at-risk population and reduces p-MRR, confirming that
the \textbf{dynamic anchor $\tau$ alone} provides optimal
discriminative power.

\subsection{Training-Set vs.\ Test-Set Derivation}
\label{sec:train-vs-test}

A critical concern in parameter derivation is \textbf{test-set leakage}:
if parameters are tuned on the test set, performance numbers are
over-optimistic and non-reproducible.  We address this by deriving all
parameters from training-set statistics alone:
%
\begin{itemize}
  \item \textbf{V5 (training-set derivation)}:
    $\alpha = 1.0, \beta = 1.926, \delta = 0.0$
    $\rightarrow$ p-MRR = 0.2360 on FollowIR test set.
  \item \textbf{V4 (test-set derivation, for comparison)}:
    $\alpha = 1.0, \beta = 1.29, \delta = 0.0$
    $\rightarrow$ p-MRR = 0.2243 on FollowIR test set.
\end{itemize}
%
Both training-set and test-set derivations independently yield
$\alpha \approx 1.0$, validating SAP as a genuine statistical
regularity of the embedding space rather than an artifact of
overfitting.  The slight difference in $\beta$ (1.926 vs.\ 1.29) is
expected: the training set captures a broader range of semantic
variation, demanding a higher Information Premium to cover the full
space of instructional constraints.  Notably, V5's higher p-MRR
(0.2360 vs.\ 0.2243) demonstrates that training-set derivation can
actually \emph{outperform} test-set grid search --- eliminating
overfitting risk while achieving better results.

This encoder-agnostic approach ensures that \method{} can be applied
to any dense retriever without manual recalibration: the parameters
are computed from the retriever's own score distributions, not from
external benchmarks.
% ============================================================
% 5. EXPERIMENTS
% ============================================================
\section{Experiments}
\label{sec:experiments}

\subsection{FollowIR Benchmark Background}
\label{sec:followir-background}

Table~\ref{tab:followir-background} reports the performance of
existing retrieval methods on the FollowIR benchmark
\cite{followir2024}.  Cross encoders (MonoT5-3B, Mistral-7B-instruct,
FollowIR-7B) and bi-encoder dense models (BM25, TART-Contriever,
Instructor XL, E5-Mistral, Google Gecko, RepLLaMA, Promptriever) are
evaluated on three TREC collections.  Most methods exhibit near-zero
or negative p-MRR, indicating poor instruction sensitivity.
Promptriever, which is fine-tuned on instruction-annotated data,
achieves the best overall score but requires supervised training.

\begin{table*}[t]
\centering
\small
\begin{tabular}{llcccccccc}
\toprule
& & \multicolumn{2}{c}{Robust04} & \multicolumn{2}{c}{News21} & \multicolumn{2}{c}{Core17} & \multicolumn{2}{c}{Average} \\
\cmidrule(lr){3-4} \cmidrule(lr){5-6} \cmidrule(lr){7-8} \cmidrule(lr){9-10}
& \textbf{Model} & MAP & p-MRR & nDCG & p-MRR & MAP & p-MRR & Score & p-MRR \\
\midrule
\multirow{3}{*}{\rotatebox[origin=c]{90}{\parbox{1.3cm}{\centering Cross\\Encoders}}}
& MonoT5-3B & 27.3 & +4.0 & 16.5 & +1.8 & 18.2 & +1.8 & 20.7 & +2.5 \\
& Mistral-7B-instruct & 23.2 & +12.6 & 27.2 & +4.8 & 19.7 & +13.0 & 23.4 & +10.1 \\
& FollowIR-7B & 24.8 & +13.7 & 29.6 & +6.3 & 20.0 & +16.5 & 24.8 & +12.2 \\
\midrule
\multirow{8}{*}{\rotatebox[origin=c]{90}{\parbox{2.2cm}{\centering Bi\\Encoders}}}
& BM25 & 12.1 & -3.1 & 19.3 & -2.1 & 8.1 & -1.1 & 13.2 & -2.1 \\
& TART-Contriever & 14.3 & -9.0 & 21.8 & -3.0 & 13.3 & -3.0 & 16.5 & -5.0 \\
& Instructor XL & 19.7 & -8.1 & 26.1 & -0.9 & 16.8 & 0.7 & 20.9 & -2.8 \\
& E5-Mistral & 23.1 & -9.6 & 27.8 & -0.9 & 18.3 & +0.1 & 23.1 & -3.5 \\
& Google Gecko & 23.3 & -2.4 & 29.5 & +3.9 & 23.2 & +5.4 & 25.3 & +2.3 \\
& RepLLaMA & 24.0 & -8.9 & 24.5 & -1.8 & 20.6 & +1.3 & 23.0 & -3.1 \\
& Promptriever & 28.3 & +11.7 & 28.5 & +6.4 & 21.6 & +15.4 & 26.1 & +11.2 \\
& DeIR-Dual & \textbf{28.6} & +8.3 & \textbf{32.3} & \textbf{+18.7} & \textbf{26.0} & \textbf{+17.6} & \textbf{29.0} & \textbf{+14.9} \\
\bottomrule
\end{tabular}
\caption{Performance of existing retrieval methods and DeIR-Dual on the FollowIR
  benchmark \cite{followir2024}.  All metrics are percentages
  ($\times 100$).  DeIR-Dual achieves the highest News21 nDCG, Core17 MAP,
  average score, and average p-MRR without any fine-tuning.}
\label{tab:followir-background}
\end{table*}

\subsection{Main Results: FollowIR Benchmark}
\label{sec:main-results}

Table~\ref{tab:followir-main} compares \method{} with six query
rewriting baselines on the FollowIR benchmark.  \method{} achieves
the highest mean p-MRR (14.9) and Score (29.0), outperforming
all baselines by 10.6$\times$ in p-MRR (vs.\ HyDE) and 50.6$\times$
(vs.\ DeepRetrieval), while also achieving the highest retrieval
quality.

\begin{table*}[t]
\centering
\small
\begin{tabular}{l|cc|cc|cc|cc|c}
\toprule
& \multicolumn{2}{c|}{Core17} & \multicolumn{2}{c|}{Robust04} & \multicolumn{2}{c|}{News21} & \multicolumn{2}{c|}{Average} & \\
\cmidrule(lr){2-3} \cmidrule(lr){4-5} \cmidrule(lr){6-7} \cmidrule(lr){8-9}
\textbf{Method} & MAP & p-MRR & MAP & p-MRR & nDCG & p-MRR & Score & p-MRR & Type \\
\midrule
DeepRetrieval & 21.5 & +6.6 & 25.9 & -5.3 & 23.9 & -2.2 & 23.8 & -0.3 & RL query gen. \\
HyDE & 23.6 & +6.5 & 24.9 & -2.9 & 21.4 & +0.7 & 23.4 & +1.4 & Hypo.\ docs \\
Query2Doc & 25.9 & +8.0 & 27.9 & -7.5 & 24.9 & -3.8 & 26.2 & -1.1 & Pseudo-doc \\
RAG-Fusion & 21.9 & +5.4 & 21.1 & -8.1 & 21.0 & +1.8 & 21.3 & -0.3 & Multi-query \\
RAG-QR & 22.3 & +2.3 & 27.6 & -10.3 & 22.9 & -2.1 & 24.3 & -3.4 & T5 rewriter \\
\midrule
DeIR-Dual & \textbf{26.0} & \textbf{+17.6} & \textbf{28.6} & \textbf{+8.3} & \textbf{32.3} & \textbf{+18.7} & \textbf{29.0} & \textbf{+14.9} & Dual-track \\
\bottomrule
\end{tabular}
\caption{Comparison of \method{} with six query rewriting methods
  on the FollowIR benchmark.  \method{} achieves the highest mean
  p-MRR (14.9) and Score (29.0), outperforming all baselines
  by 10.6$\times$ in p-MRR (vs.\ HyDE) and 50.6$\times$ (vs.\
  DeepRetrieval).  All metrics are percentages ($\times 100$).}
\label{tab:followir-main}
\end{table*}

\subsection{Cross-Dataset Validation}
\label{sec:cross-dataset}

\todo{NegConstraint (text retrieval with negation signals);
  COCO-Neg (multimodal text-to-image negation). Both show
  significant improvement (+4.3\% nDCG@10 on NegConstraint).}

\subsection{BEIR Generalization}
\label{sec:beir-generalization}

\todo{NQ (+1.1\% nDCG@10), HotpotQA (-0.5\%, within variance),
  MS MARCO (+2.0\% nDCG@10). CONSERVATIVE prompting avoids
  pseudo-negation disaster.}

\subsection{Ablation Study}
\label{sec:ablation}

\todo{Ablate each V2 component: dynamic $\tau$, safety gate,
  gap weighting, Softplus penalty, $\qminus$ contribution.}

\subsection{Encoder-Agnostic Analysis}
\label{sec:encoder-agnostic}

\todo{EAPS strategy: RepLLaMA (0.08\% at-risk) vs Mistral (62.9\%
  at-risk). Retrieval-simulated distractor sampling with top-k=1000.}

% ============================================================
% 6. ANALYSIS: SEMANTIC ENTANGLEMENT
% ============================================================
\section{Analysis: Semantic Entanglement}
\label{sec:analysis}

\subsection{Definition and Variants}
\label{sec:entanglement-definition}

\todo{Structural entanglement (ComLQ: negation clauses share semantics
  with positive intent). Generative entanglement (NQ: LLM generates
  pseudo-negations overlapping with $\qbase$).}

\subsection{Diagnostic Analysis}
\label{sec:entanglement-diagnosis}

\todo{$\cos(\mathbf{h}_{\text{base}}, \mathbf{h}_{\text{neg}})$ distribution: NegConstraint (orthogonal)
  vs ComLQ (0.705 mean) vs NQ (0.625 mean). $S_{\text{neg}}$ higher
  on relevant documents $\rightarrow$ double whammy effect.}

\begin{figure}[t]
  \centering
  % \includegraphics[width=\columnwidth]{figures/semantic_entanglement.pdf}
  \fbox{\parbox{\columnwidth}{\centering\vspace{3cm}%
    \textbf{Figure 3: Semantic Entanglement Distribution}\\
    \small (placeholder: to be replaced with violin/box plot)\vspace{1cm}}}
  \caption{Distribution of $\cos(\mathbf{h}_{\text{base}}, \mathbf{h}_{\text{neg}})$ across three
    scenarios. NegConstraint shows orthogonality (penalty effective),
    while ComLQ (structural entanglement) and NQ (generative
    entanglement) exhibit significant overlap, causing penalty failure.}
  \label{fig:semantic-entanglement}
\end{figure}

% ============================================================
% 7. CONCLUSION AND LIMITATIONS
% ============================================================
\section{Conclusion and Limitations}
\label{sec:conclusion}

\todo{Summary: three contributions. Limitations: semantic entanglement,
  LLM dependency, high at-risk encoder sensitivity. Future directions.}

% ============================================================
% REFERENCES
% ============================================================
\bibliographystyle{plain}
\bibliography{../refs/references}

\end{document}
